% ---- pg-303 （书页 291） ----

\mvsub{回归参数的似然比检验}

回归分析的部分内容涉及评估特定预测变量对响应变量的影响. 有一个人们感兴趣的零假设是说某些 $z_i$ 对响应 $Y$ 没有影响. 可将这些预测变量表示为 $z_{q+1},z_{q+2},\cdots,z_r$. 对 $Y$ 无影响这种说法表示成统计假设
\phantomsection\label{eq:7-12}
\begin{equation}
H_0:\beta_{q+1}=\beta_{q+2}=\cdots=\beta_r=0\qquad H_0:\boldsymbol{\beta}_{(2)}=\mathbf{0}
\tag{7-12}
\end{equation}
其中 $\boldsymbol{\beta}_{(2)}'=[\beta_{q+1},\beta_{q+2},\cdots,\beta_r]'$. 令
\[
\mathbf{Z}=\left[\begin{array}{c:c}
\underset{n\times(q+1)}{\mathbf{Z}_1}&\underset{n\times(r-q)}{\mathbf{Z}_2}
\end{array}\right],\qquad
\boldsymbol{\beta}=\left[\begin{array}{c}
\underset{(q+1)\times 1}{\boldsymbol{\beta}_{(1)}}\\\hdashline
\underset{(r-q)\times 1}{\boldsymbol{\beta}_{(2)}}
\end{array}\right]
\]
我们可将一般线性模型表示为
\[
\mathbf{Y}=\mathbf{Z}\boldsymbol{\beta}+\boldsymbol{\epsilon}
=\left[\begin{array}{c:c}\mathbf{Z}_1&\mathbf{Z}_2\end{array}\right]
\begin{bmatrix}\boldsymbol{\beta}_{(1)}\\\boldsymbol{\beta}_{(2)}\end{bmatrix}
+\boldsymbol{\epsilon}
=\mathbf{Z}_1\boldsymbol{\beta}_{(1)}+\mathbf{Z}_2\boldsymbol{\beta}_{(2)}+\boldsymbol{\epsilon}
\]
在假设 $H_0:\boldsymbol{\beta}_{(2)}=\mathbf{0}$ 下, $\mathbf{Y}=\mathbf{Z}_1\boldsymbol{\beta}_{(1)}+\boldsymbol{\epsilon}$. $H_0$ 的似然比检验的依据是
\phantomsection\label{eq:7-13}
\begin{equation}
\begin{aligned}
\text{额外的平方和}&=\text{SS}_{\text{残差}}(\mathbf{Z}_1)-\text{SS}_{\text{残差}}(\mathbf{Z})\\
&=(\mathbf{y}-\mathbf{Z}_1\hat{\boldsymbol{\beta}}_{(1)})'(\mathbf{y}-\mathbf{Z}_1\hat{\boldsymbol{\beta}}_{(1)})-(\mathbf{y}-\mathbf{Z}\hat{\boldsymbol{\beta}})'(\mathbf{y}-\mathbf{Z}\hat{\boldsymbol{\beta}})
\end{aligned}
\tag{7-13}
\end{equation}
其中 $\hat{\boldsymbol{\beta}}_{(1)}=(\mathbf{Z}_1'\mathbf{Z}_1)^{-1}\mathbf{Z}_1'\mathbf{y}$.

\mvthm{结论 7.6}{设 $\mathbf{Z}$ 有满秩 $r+1$, 且 $\boldsymbol{\epsilon}$ 服从 $N_n(\mathbf{0},\sigma^2\mathbf{I})$ 分布. 假设 $H_0:\boldsymbol{\beta}_{(2)}=\mathbf{0}$ 的似然比检验等价于 $H_0$ 按\hyperref[eq:7-13]{式(7-13)} 中额外平方和以及 $s^2=(\mathbf{y}-\mathbf{Z}\hat{\boldsymbol{\beta}})'(\mathbf{y}-\mathbf{Z}\hat{\boldsymbol{\beta}})/(n-r-1)$ 所作的检验. 特别是似然比检验对 $H_0$ 的拒绝域为}
\[
\frac{(\text{SS}_{\text{残差}}(\mathbf{Z}_1)-\text{SS}_{\text{残差}}(\mathbf{Z}))/(r-q)}{s^2}>F_{r-q,n-r-1}(\alpha)
\]
其中 $F_{r-q,n-r-1}(\alpha)$ 为自由度为 $r-q$ 和 $n-r-1$ 的 $F$ 分布的上 $100\alpha$ 百分位数.

证明\quad 给定数据和正态性假定, 参数 $\boldsymbol{\beta}$ 和 $\sigma^2$ 的似然函数为
\[
L(\boldsymbol{\beta},\sigma^2)=\frac{1}{(2\pi)^{n/2}\sigma^n}e^{-(\mathbf{y}-\mathbf{Z}\boldsymbol{\beta})'(\mathbf{y}-\mathbf{Z}\boldsymbol{\beta})/2\sigma^2}
\leq\frac{1}{(2\pi)^{n/2}\hat{\sigma}^n}e^{-n/2}
\]
最大值在 $\hat{\boldsymbol{\beta}}=(\mathbf{Z}'\mathbf{Z})^{-1}\mathbf{Z}'\mathbf{y}$, 及 $\hat{\sigma}^2=(\mathbf{y}-\mathbf{Z}\hat{\boldsymbol{\beta}})'(\mathbf{y}-\mathbf{Z}\hat{\boldsymbol{\beta}})/n$ 时达到. 而在零假设的约束下, $\mathbf{Y}=\mathbf{Z}_1\boldsymbol{\beta}_{(1)}+\boldsymbol{\epsilon}$ 且
\[
\max_{\boldsymbol{\beta}_{(1)},\sigma^2}L(\boldsymbol{\beta}_{(1)},\sigma^2)=\frac{1}{(2\pi)^{n/2}\hat{\sigma}_1^n}e^{-n/2}
\]
其中最大值在 $\hat{\boldsymbol{\beta}}_{(1)}=(\mathbf{Z}_1'\mathbf{Z}_1)^{-1}\mathbf{Z}_1'\mathbf{y}$ 时出现, 此外,
\[
\hat{\sigma}_1^2=(\mathbf{y}-\mathbf{Z}_1\hat{\boldsymbol{\beta}}_{(1)})'(\mathbf{y}-\mathbf{Z}_1\hat{\boldsymbol{\beta}}_{(1)})/n
\]
似然比为
\[
\frac{\max\limits_{\boldsymbol{\beta}_{(1)},\sigma^2}L(\boldsymbol{\beta}_{(1)},\sigma^2)}{\max\limits_{\boldsymbol{\beta},\sigma^2}L(\boldsymbol{\beta},\sigma^2)}
=\left(\frac{\hat{\sigma}_1^2}{\hat{\sigma}^2}\right)^{-n/2}
=\left(\frac{\hat{\sigma}^2+\hat{\sigma}_1^2-\hat{\sigma}^2}{\hat{\sigma}^2}\right)^{-n/2}
=\left(1+\frac{\hat{\sigma}_1^2-\hat{\sigma}^2}{\hat{\sigma}^2}\right)^{-n/2}
\]
当似然比值很小时拒绝 $H_0:\boldsymbol{\beta}_{(2)}=\mathbf{0}$, 这等价于当 $(\hat{\sigma}_1^2-\hat{\sigma}^2)/\hat{\sigma}^2$ 值很大时拒绝 $H_0$, 或当
\[
\frac{n(\hat{\sigma}_1^2-\hat{\sigma}^2)/(r-q)}{n\hat{\sigma}^2/(n-r-1)}
=\frac{(\text{SS}_{\text{残差}}(\mathbf{Z}_1)-\text{SS}_{\text{残差}}(\mathbf{Z}))/(r-q)}{s^2}=F
\]
